\section*{Введение}
\addcontentsline{toc}{section}{Введение}
% \addcontentsline{toc}{section}{\protect\numberline{}Введение}
\label{sec:Chapter0} \index{Chapter0}

 	
В качестве модельной задачи предлагается задача для трехмерного гиперболического уравнения в области, представляющей из себя прямоугольный параллелепипед.
Индивидуальные варианты заданий отличаются типом граничных условий.
Данное уравнение часто применяется в теории тепло и массопереноса, гидро и аэромеханике, электростатике.
Поэтому поиск решения данный задачи в различных областях является весьма актуальным.


Задание необходимо выполнить на следующих ПВС Московского университета:
\begin{enumerate}
    \item IBM Polus \cite{Polus}.
\end{enumerate}
